% =====================================================================
%  Lecture 04 -- Mobility and Handover                    (English edition)
% =====================================================================
\documentclass[11pt,a4paper]{article}
\providecommand\ACNroot{../..}
\input{\ACNroot/templates/lecture-en}

\ACNset{number}{04}
\ACNset{title}{Mobility and Handover}
\ACNset{shorttitle}{Mobility and Handover}
\ACNset{titleother}{تحرک و تحویل}
\ACNset{part}{Part I -- 2G/3G foundations (Lectures 1--4)}
\ACNset{time}{7--9 hours (3.5 h reading, 4--5 h exercises)}
\ACNset{level}{Intermediate}
\ACNset{prereq}{Lectures 1--3; log-distance path loss; Gaussian distribution and the $Q$-function}
\ACNset{gen}{4}

\begin{document}
\ACNcover

% ---------------------------------------------------------------------
\section{Learning objectives}
After studying this lecture you should be able to:
\begin{objectives}
  \item \textbf{Distinguish} idle-mode mobility (cell reselection, location registration) from
        connected-mode mobility (handover) and explain who decides in each case.
  \item \textbf{Compare} location areas, routing areas, tracking-area lists and RAN notification
        areas, and \textbf{optimise} the size of a registration area with a cost model.
  \item \textbf{Classify} handovers (intra-/inter-frequency, inter-RAT, hard/soft) and state the
        measurement quantities RSRP and RSRQ.
  \item \textbf{Apply} the LTE A3 entering condition with hysteresis, offsets and time-to-trigger,
        and \textbf{compute} where a handover is triggered.
  \item \textbf{Trace} an X2-based LTE handover message by message, including data forwarding and
        path switch.
  \item \textbf{Diagnose} too-late, too-early, wrong-cell and ping-pong handovers and relate them
        to parameter settings.
\end{objectives}

% ---------------------------------------------------------------------
\section{Prerequisite knowledge}
You need the location-area concept of Lectures~1--2, the soft handover of WCDMA (Lecture~3),
and the log-distance path-loss model with log-normal shadowing:
$\PL(d) = \PL(d_0) + 10n\log_{10}(d/d_0) + X_\sigma$, where $X_\sigma$ is a zero-mean Gaussian
variable (in dB) with standard deviation $\sigma$. The tail function is
$Q(x) = \Pr\{Z > x\}$ for a standard Gaussian $Z$.

\paragraph{Key terms} The vocabulary of this lecture with its Persian equivalents:
\begin{keyterms}
Mobility management & \fa{مدیریت تحرک} & Keeping a moving user reachable and connected. \\
Cell reselection & \fa{انتخاب مجدد سلول} & UE-controlled change of camping cell in idle mode. \\
Handover (HO) & \fa{تحویل} & Network-controlled transfer of an active connection to another cell. \\
Tracking area (TA) & \fa{ناحیهٔ ردیابی} & LTE/5G registration area; the paging scope. \\
RSRP & \fa{توان دریافتی سیگنال مرجع} & Average power of one cell-specific reference-signal resource element. \\
RSRQ & \fa{کیفیت دریافتی سیگنال مرجع} & $N\cdot\mathrm{RSRP}/\mathrm{RSSI}$: a load- and interference-aware quality. \\
Hysteresis & \fa{پسماند} & Margin that prevents repeated triggering around a threshold. \\
Time-to-trigger (TTT) & \fa{زمان تا راه‌اندازی} & Duration a condition must hold before a report is sent. \\
Ping-pong & \fa{پینگ‌پنگ} & Unnecessary handover followed quickly by a return. \\
Radio link failure (RLF) & \fa{شکست پیوند رادیویی} & Loss of the connection detected by timers (e.g.\ T310). \\
\end{keyterms}

% ---------------------------------------------------------------------
\section{Introduction and motivation}
A fixed telephone is always at the same socket; a mobile user is not. Mobility management
answers two different questions:
\begin{enumerate}[label=(\alph*)]
  \item \textbf{Reachability} -- when the user is idle, where should the network look for it
        when a call or a packet arrives?
  \item \textbf{Session continuity} -- when the user is active, how can the connection follow
        the user from cell to cell without dropping or interrupting it?
\end{enumerate}
Both questions involve a trade-off between \emph{signalling cost} and \emph{performance}. If
the network knew the exact cell of every idle user, every cell change would cost an update; if
it knew nothing, every incoming call would have to be paged everywhere. If handovers are
triggered eagerly, users bounce between cells (ping-pong); if they are triggered late, links
fail. This lecture develops both sides across 2G to 5G and gives each trade-off a quantitative
form.

% ---------------------------------------------------------------------
\section{Fundamental concepts}

\subsection{Idle mode versus connected mode}
In \textbf{idle mode} the UE has no radio connection. It \emph{camps} on a cell, listens to its
broadcast and paging channels, and chooses by itself when to move to a better cell (cell
reselection). The network only knows the UE's registration area. In \textbf{connected mode}
the UE has a dedicated radio connection; the network decides the serving cell, based on
measurement reports from the UE (network-controlled, mobile-assisted handover).

\subsection{Registration areas across generations}
\begin{itemize}
  \item \textbf{Location area (LA)} -- GSM/UMTS CS domain; updated towards the VLR.
  \item \textbf{Routing area (RA)} -- GPRS/UMTS PS domain; a subset of an LA, updated towards
        the SGSN. Two parallel procedures were needed because the two domains were separate.
  \item \textbf{Tracking area (TA)} -- LTE and 5G. The MME (AMF in 5G) gives the UE a
        \emph{list} of TAs; the UE updates only when it enters a TA outside the list. Lists
        can be tailored per UE, which removes the ping-pong of updates at a fixed border.
  \item \textbf{RAN notification area (RNA)} -- 5G RRC\_INACTIVE state: the core considers the UE
        connected, the RAN keeps its context and pages it within a RAN-configured area.
\end{itemize}

\begin{acnfigure}
  \resizebox{\linewidth}{!}{\acnfig{lecture-04/area-hierarchy}}
  \caption{Registration areas of successive generations.}
  \label{fig:areas}
\end{acnfigure}

\subsection{Measurement quantities}
LTE defines two basic measurements on the cell-specific reference signals
\cite{ts36300}:
\begin{definitionbox}{RSRP and RSRQ}
\emph{RSRP} is the linear average of the power of the resource elements that carry
cell-specific reference signals within the measured bandwidth (range about $-140$ to
$-44$\,dBm). \emph{RSRQ} is
\begin{equation}\label{eq:rsrq}
  \RSRQ = \frac{N\cdot\RSRP}{\RSSI},
\end{equation}
where $N$ is the number of resource blocks in the RSSI measurement bandwidth and RSSI is the total
received power (wanted signal, interference and noise) over that bandwidth.
\end{definitionbox}
RSRP measures coverage; RSRQ also reflects load and interference, because RSSI grows when
neighbours are busy. NR uses the analogous SS-RSRP and SS-RSRQ on synchronisation signal blocks.

% ---------------------------------------------------------------------
\section{Detailed technical explanation}

\subsection{Cell reselection in idle mode}
In LTE a cell is \emph{suitable} if it satisfies the S-criterion \cite{ts36304}:
\begin{equation}\label{eq:scrit}
  S_{\mathrm{rxlev}} = Q_{\mathrm{rxlevmeas}} - (Q_{\mathrm{rxlevmin}} + Q_{\mathrm{rxlevminoffset}})
  - P_{\mathrm{compensation}} > 0 ,
\end{equation}
where $Q_{\mathrm{rxlevmeas}}$ is the measured RSRP, $Q_{\mathrm{rxlevmin}}$ the broadcast minimum
level and $P_{\mathrm{compensation}}$ a penalty for UEs whose maximum power is below the level
assumed by the operator. Among suitable intra-frequency cells the UE ranks the serving and
neighbour cells with
\begin{equation}\label{eq:rcrit}
  R_\mathrm{s} = Q_{\mathrm{meas,s}} + Q_{\mathrm{hyst}}, \qquad
  R_\mathrm{n} = Q_{\mathrm{meas,n}} - Q_{\mathrm{offset}},
\end{equation}
and reselects a neighbour only if it is ranked best for a time $T_{\mathrm{reselection}}$.
Inter-frequency and inter-RAT reselection use absolute \emph{priorities}: a UE always moves to a
higher-priority layer that is good enough, and leaves to a lower-priority layer only when its
current layer becomes poor. Priorities let operators steer idle users towards, for example, a
capacity layer or a 5G carrier.

\subsection{Handover classification}
\begin{itemize}
  \item By \textbf{frequency}: \emph{intra-frequency} (same carrier; measurable without gaps),
        \emph{inter-frequency} (another carrier; the UE needs measurement gaps unless it has a
        second receiver) and \emph{inter-RAT} (another technology, e.g.\ LTE to UMTS).
  \item By \textbf{connection continuity}: \emph{hard} handover (break-before-make; GSM, LTE, NR)
        and \emph{soft} handover (make-before-break; WCDMA, where several cells transmit the same
        data simultaneously).
  \item By \textbf{network scope}: within one controller (intra-BSC, intra-eNB), between base
        stations (X2/Xn handover) or involving the core (inter-MSC in GSM, S1-based in LTE, N2-based
        in 5G).
\end{itemize}

\paragraph{Handover in GSM} The BTS measures the uplink and the MS reports the six strongest
neighbours every 480\,ms on the SACCH. The BSC decides; inter-BSC handovers involve the MSC, and
inter-MSC handovers keep the first MSC as an \emph{anchor} for the call \cite{ts23009}.

\paragraph{Soft handover in WCDMA} The UE maintains an \emph{active set} of cells; events add a
cell that comes within a reporting range of the best cell, remove a cell that falls out of it,
or replace the weakest member. The cost is extra downlink power and Iub capacity, typically
30--40\% of users being in soft handover.

\paragraph{Handover in LTE} LTE removed the RNC and soft handover. The eNB decides alone, based
on measurement \emph{events} configured through RRC \cite{ts36331}:
\begin{table}[htbp]
\centering\small
\caption{LTE intra-RAT measurement events (thresholds and offsets are configured by RRC).}\label{tab:events}
\rowcolors{2}{acnmist}{white}
\begin{tabularx}{\linewidth}{@{}l Y Y@{}}
\toprule
\rowcolor{white}Event & Entering condition (simplified) & Typical use \\
\midrule
A1 & Serving becomes better than a threshold & Stop inter-frequency measurements (remove gaps) \\
A2 & Serving becomes worse than a threshold & Start inter-frequency/inter-RAT measurements \\
A3 & Neighbour becomes offset better than serving (PCell) & Intra-frequency handover \\
A4 & Neighbour becomes better than a threshold & Load-based or coverage-layer handover \\
A5 & Serving worse than threshold~1 and neighbour better than threshold~2 & Inter-frequency handover \\
A6 & Neighbour becomes offset better than a secondary cell & Carrier-aggregation management \\
B1/B2 & Inter-RAT neighbour better than a threshold (and serving worse, B2) & LTE to UMTS/GSM, LTE to NR \\
\bottomrule
\end{tabularx}
\end{table}

\subsection{Handover in 5G NR}
NR keeps the LTE principles (eNB$\to$gNB, X2$\to$Xn, S1$\to$N2) and adds features aimed at the
interruption and the failures of hard handover \cite{ts38300}:
\begin{itemize}
  \item \textbf{Conditional handover (CHO)}: the network prepares one or more target cells in
        advance and sends the UE a handover command with an execution condition; the UE executes
        it when the condition is met, avoiding the race in which the command arrives over a
        failing link.
  \item \textbf{Dual active protocol stack (DAPS)}: the UE keeps receiving from the source while
        connecting to the target, approaching zero interruption.
  \item \textbf{Beam management} inside a cell: many beam changes are handled at layers 1/2
        without RRC handover.
\end{itemize}

\subsection{Handover failures}
A UE declares \emph{radio link failure} (RLF) when, after $N310$ consecutive out-of-sync
indications, timer $T310$ expires without recovery; it then attempts RRC re-establishment.
Self-organising networks classify failures to tune parameters (mobility robustness
optimisation, MRO):
\begin{itemize}
  \item \textbf{Too late} -- RLF in the source cell before (or during) the handover; the UE
        re-establishes in the target. Remedy: earlier triggering (lower offset/TTT).
  \item \textbf{Too early} -- RLF shortly after a successful handover; the UE re-establishes in the
        \emph{source}. Remedy: later triggering.
  \item \textbf{Wrong cell} -- RLF shortly after handover; re-establishment in a third cell.
        Remedy: adjust the cell-individual offsets of the two neighbours.
  \item \textbf{Ping-pong} -- a successful handover followed by a return to the source within a
        short time; wastes signalling and causes interruptions. Remedy: more hysteresis or TTT.
\end{itemize}
Too-late and ping-pong pull the parameters in opposite directions: this is the central
tuning dilemma of handover.

% ---------------------------------------------------------------------
\section{Mathematical models}

\subsection{The A3 condition and the handover point}\label{sec:a3}
\paragraph{Engineering problem} Where along a road between two cells does a handover actually
happen, and how do the parameters move that point?
\paragraph{The A3 entering condition} In LTE \cite{ts36331} the condition, with all quantities in
dB or dBm, is
\begin{equation}\label{eq:a3}
  M_\mathrm{n} + O_{\mathrm{fn}} + O_{\mathrm{cn}} - \Hys \;>\; M_\mathrm{p} + O_{\mathrm{fp}} + O_{\mathrm{cp}} + \Off ,
\end{equation}
where $M_\mathrm{n}$, $M_\mathrm{p}$ are the (layer-3 filtered) measurements of the neighbour and
the primary serving cell, $O_{\mathrm{f}\cdot}$ frequency-specific offsets, $O_{\mathrm{c}\cdot}$
cell-individual offsets, $\Hys$ the hysteresis and $\Off$ the event offset. The condition must
hold continuously for $\TTT$ before the UE reports. With zero frequency and cell offsets it reduces
to $M_\mathrm{n} - M_\mathrm{p} > H$ with $H = \Off + \Hys$.

\paragraph{System model} Two eNBs at distance $D$ with equal transmit power; the UE moves along the
line joining them; deterministic path loss with exponent $n$ (no shadowing). At distance $x$ from
the serving eNB,
\begin{equation}
  M_\mathrm{n} - M_\mathrm{p} = 10n\log_{10}\frac{x}{D-x}.
\end{equation}
\paragraph{Derivation} Setting $10n\log_{10}\bigl(x/(D-x)\bigr) = H$ and solving,
\begin{equation}\label{eq:hopoint}
  x_{\mathrm{HO}} = D\,\frac{r}{1+r}, \qquad r = 10^{H/(10n)} .
\end{equation}
The UE then travels a further $v\,\TTT$ before reporting.
\paragraph{Interpretation} For $D = \qty{1000}{\metre}$, $n=3.5$ and $H = \qty{3}{\dB}$: $r=1.218$
and $x_{\mathrm{HO}} = \qty{549}{\metre}$, i.e.\ 49\,m past the cell border; $H = \qty{6}{\dB}$ moves
it to 597\,m. At \qty{30}{\metre\per\second} a TTT of 320\,ms adds another 9.6\,m
(\cref{fig:a3}). Every decibel of margin is paid for with signal quality on the old cell, where
the link is already the weakest -- the origin of too-late failures at high speed.

\begin{acnfigure}
  \acnfig{lecture-04/a3-trigger}
  \caption{A3 triggering: the neighbour must exceed the serving cell by $\Off + \Hys = 3$\,dB
  (dotted) for $\TTT = 320$\,ms before the measurement report is sent.}
  \label{fig:a3}
\end{acnfigure}

\subsection{Shadowing and ping-pong}\label{sec:pingpong}
\paragraph{Model} Add independent log-normal shadowing $X_\mathrm{s}, X_\mathrm{n}\sim\mathcal{N}(0,\sigma^2)$
(dB) to both cells. At the cell border the mean difference is zero, and
$M_\mathrm{n} - M_\mathrm{p} \sim \mathcal{N}(0, 2\sigma^2)$.
\paragraph{Result} The probability that a single sample at the border exceeds the margin $H$ is
\begin{equation}\label{eq:pingpong}
  P_{\mathrm{trig}} = Q\!\left(\frac{H}{\sqrt{2}\,\sigma}\right).
\end{equation}
With $\sigma = \qty{8}{\dB}$ and $H = \qty{3}{\dB}$, $P_{\mathrm{trig}} = Q(0.265) = 0.40$; even with
$H=\qty{6}{\dB}$ it is still $0.30$. Hysteresis alone cannot suppress spurious triggers caused by
large-scale fading; TTT and layer-3 filtering are needed. The layer-3 filter
\begin{equation}\label{eq:l3filter}
  F_k = (1-a)\,F_{k-1} + a\,M_k, \qquad a = 2^{-K/4},
\end{equation}
with filter coefficient $K$, reduces the variance of \emph{independent} fluctuations (fast
fading, measurement error) by the factor $a/(2-a)$ -- one third for $K=4$ -- but it cannot remove
shadowing, which is correlated over tens of metres. The model is pessimistic in one respect
(it ignores the correlation of shadowing between the two cells) and optimistic in another (it
looks at one point only); simulation along trajectories is the usual next step.

\subsection{Sizing a registration area}\label{sec:lasize}
\paragraph{Problem} How many cells should a location or tracking area contain?
\paragraph{Model and assumptions} Square cells of side $a$; a registration area of $N$ cells
forms a square of side $a\sqrt{N}$ and perimeter $L = 4a\sqrt{N}$. Users are uniformly distributed
and move in uniformly random directions at mean speed $v$ (the \emph{fluid-flow} model
\cite{akyildiz1999mobility}): the rate at which users of an area $S$ cross its perimeter $L$ is
$\rho vL/\pi$, where $\rho$ is the user density. Each user receives $\lambda_\mathrm{c}$ calls per
unit time, each paged in all $N$ cells. One update costs $c_\mathrm{u}$ and paging one cell costs
$c_\mathrm{p}$.
\paragraph{Derivation} Per user, the update rate is $\rho vL/(\pi\rho S) = vL/(\pi S)$ with
$S = Na^2$, so
\begin{equation}\label{eq:cost}
  C(N) = c_\mathrm{u}\,\frac{4v}{\pi a\sqrt{N}} + c_\mathrm{p}\,\lambda_\mathrm{c}\,N .
\end{equation}
The first term decreases and the second increases with $N$. Setting $\mathrm{d}C/\mathrm{d}N = 0$,
\begin{equation}\label{eq:nopt}
  -\frac{2c_\mathrm{u}v}{\pi a}N^{-3/2} + c_\mathrm{p}\lambda_\mathrm{c} = 0
  \quad\Longrightarrow\quad
  N^\ast = \left(\frac{2\,c_\mathrm{u}\,v}{\pi\,a\,c_\mathrm{p}\,\lambda_\mathrm{c}}\right)^{2/3}.
\end{equation}
\paragraph{Interpretation} Fast users ($v$ large) favour large areas; users who receive many calls
($\lambda_\mathrm{c}$ large) favour small ones. Because the optimum depends on the user, LTE's
per-UE tracking-area lists can approach it user by user, which a fixed LA plan cannot.

\begin{acnfigure}
  \acnfig{lecture-04/la-cost}
  \caption{Update, paging and total signalling cost of \cref{eq:cost} for
  $c_\mathrm{u} = 1$, $c_\mathrm{p} = 0.1$, $v = \qty{10}{\metre\per\second}$, $a = \qty{500}{\metre}$,
  $\lambda_\mathrm{c} = 2$ calls per hour (\cref{ex:la}).}
  \label{fig:lacost}
\end{acnfigure}

% ---------------------------------------------------------------------
\section{Network architecture and diagrams}
The LTE handover architecture is flat: the source and target eNBs negotiate directly over X2,
and the MME and S-GW are involved only at the end to move the downlink path (\emph{path switch}).
If no X2 exists, or if the MME or S-GW must change, the same steps run through the MME over S1
(\emph{S1-based handover}). In 5G the roles map to gNB, Xn, AMF/SMF and UPF (Lecture~20). GSM
and UMTS involve their controllers (BSC/RNC) and, across controller borders, the MSC/SGSN.

% ---------------------------------------------------------------------
\section{Step-by-step procedures}

\begin{procedurebox}{LTE X2-based handover}\label{proc:x2}
\begin{steps}
  \item \textbf{Preparation trigger.} The UE sends a \msg{Measurement Report} after the A3 condition
        held for TTT. The source eNB decides to hand over.
  \item \textbf{Preparation.} The source sends X2AP \msg{Handover Request} with the UE context (bearers,
        security, RRC configuration). The target performs admission control, reserves resources
        (including a dedicated RACH preamble) and answers \msg{Handover Request Acknowledge}
        containing the RRC message for the UE.
  \item \textbf{Execution.} The source sends \msg{RRC Connection Reconfiguration} with mobility
        control information, stops scheduling the UE, sends \msg{SN Status Transfer} and starts
        forwarding buffered and in-transit downlink packets to the target over X2.
  \item \textbf{Access.} The UE synchronises to the target, performs contention-free random access
        (Lecture~8), receives a timing advance and an uplink grant, and sends
        \msg{RRC Connection Reconfiguration Complete}. The interruption ends.
  \item \textbf{Completion.} The target asks the MME to move the path with S1AP
        \msg{Path Switch Request}; the MME sends \msg{Modify Bearer Request} to the S-GW, which switches
        the downlink tunnel and sends \emph{end marker} packets on the old path so that the target
        knows when forwarded data is finished.
  \item \textbf{Release.} After \msg{Path Switch Request Acknowledge}, the target sends
        \msg{UE Context Release}; the source frees its resources.
\end{steps}
\end{procedurebox}

\begin{acnfigure}
  \resizebox{\linewidth}{!}{\acnfig{lecture-04/x2-handover}}
  \caption{X2-based handover in LTE (\cref{proc:x2}). Teal arrows: user-plane data forwarding and
  end marker.}
  \label{fig:x2}
\end{acnfigure}

The user-plane interruption lasts from the moment the UE leaves the source until it can be
scheduled in the target; it is of the order of a few tens of milliseconds in LTE and is the
quantity that RACH-less handover, make-before-break and DAPS reduce.

% ---------------------------------------------------------------------
\section{Worked examples}

\begin{examplebox}{Handover point and speed}\label{ex:hopoint}
\textbf{Problem.} Inter-site distance 1000\,m, $n=3.5$, $\Off+\Hys = 3$\,dB, $\TTT = 320$\,ms.
Where is the measurement report sent for a pedestrian (\qty{1.5}{\metre\per\second}) and for a car on a
motorway (\qty{120}{\kilo\metre\per\hour})?

\textbf{Solution.} From \cref{eq:hopoint}, $x_{\mathrm{HO}} = 549$\,m. The pedestrian adds
$1.5\times0.32 = 0.5$\,m; the car adds $33.3\times0.32 = 10.7$\,m, reporting at $\approx\qty{560}{\metre}$.

\textbf{Interpretation.} Then preparation and execution take further tens of milliseconds, during
which the car moves another 1--3\,m per 30\,ms. Motorway cells are therefore often given a lower
offset or TTT, and speed-dependent scaling of TTT is part of the standard.
\end{examplebox}

\begin{examplebox}{RSRQ under load}\label{ex:rsrq}
\textbf{Problem.} A UE measures RSRP $=-90$\,dBm over $N=50$ resource blocks. Compute RSRQ when the
RSSI is $-65$\,dBm (lightly loaded neighbours) and $-60$\,dBm (heavily loaded).

\textbf{Solution.} $\RSRQ = 10\log_{10}50 + \RSRP - \RSSI = 17.0 - 90 + 65 = \qty{-8.0}{\dB}$, and
$17.0 - 90 + 60 = \qty{-13.0}{\dB}$.

\textbf{Interpretation.} Same coverage, 5\,dB worse quality: RSRQ sees the load of the
neighbours, RSRP does not. Using RSRQ for A3 spreads load, at the cost of more volatile reports.
\end{examplebox}

\begin{examplebox}{Tracking-area size}\label{ex:la}
\textbf{Problem.} With the parameters of \cref{fig:lacost}, find the optimal number of cells per
area and the cost per user and hour.

\textbf{Solution.} $\lambda_\mathrm{c} = 2/3600\,\mathrm{s}^{-1}$. \Cref{eq:nopt}:
$N^\ast = \bigl(2\times1\times10/(\pi\times500\times0.1\times2/3600)\bigr)^{2/3} = (229.2)^{2/3} = 37.5$.
A $6\times6$ area ($N=36$) gives $C = 3600\,[4\times10/(\pi\times500\times6) + 0.1\times(2/3600)\times36]
= 15.28 + 7.20 = 22.5$ cost units per user-hour, against 91.9 for $N=1$ and 29.2 for $N=100$.

\textbf{Interpretation.} The optimum is flat: anything between about 20 and 60 cells costs within
10\% of the minimum. Operators can therefore choose area borders for practical reasons (along
low-traffic lines, away from busy roads) without losing much.
\end{examplebox}

% ---------------------------------------------------------------------
\section{Technology comparison}
\begin{table}[htbp]
\centering\small
\caption{Connected-mode mobility from 2G to 5G.}\label{tab:compare}
\rowcolors{2}{acnmist}{white}
\begin{tabularx}{\linewidth}{@{}l Y Y Y Y@{}}
\toprule
\rowcolor{white} & GSM & UMTS (WCDMA) & LTE & NR (5G) \\
\midrule
Type            & Hard            & Soft (same freq.) + hard & Hard                   & Hard; CHO and DAPS options \\
Decision        & BSC             & RNC                      & Source eNB             & Source gNB \\
Measurements    & RXLEV, RXQUAL   & CPICH RSCP, $E_\mathrm{c}/N_0$ & RSRP, RSRQ, SINR & SS-RSRP, SS-RSRQ, SS-SINR, CSI-RS \\
Inter-node link & A (via MSC)     & Iur                      & X2 (or S1)             & Xn (or N2) \\
Core involvement& MSC anchor      & SRNC relocation          & Path switch (MME, S-GW)& Path switch (AMF, SMF, UPF) \\
Idle area       & LA (RA for PS)  & LA + RA                  & TA list                & TA list + RNA \\
\bottomrule
\end{tabularx}
\end{table}

\paragraph{Hard versus soft} Soft handover gives diversity gain and no interruption, but needs
same-frequency neighbours, a central combining point (the RNC) and extra resources. LTE chose a
flat architecture without a central RAN node, so soft handover was impossible; it compensated
with fast hard handover, data forwarding and, later, DAPS. The trade-off --
\emph{centralised combining versus flat, fast decisions} -- reappears in 5G with the CU/DU split
(Lecture~14).

% ---------------------------------------------------------------------
\section{Key takeaways}
\begin{takeawaybox}
\begin{itemize}
  \item Idle mobility is UE-controlled (reselection, area updates); connected mobility is
        network-controlled and mobile-assisted (measurement reports, handover).
  \item Area size trades update cost against paging cost: $N^\ast = (2c_\mathrm{u}v/(\pi a c_\mathrm{p}\lambda_\mathrm{c}))^{2/3}$;
        TA lists let the network tailor it per UE.
  \item RSRP measures coverage, RSRQ $= N\,\RSRP/\RSSI$ adds load and interference.
  \item A3: $M_\mathrm{n} - M_\mathrm{p} > \Off + \Hys$ (plus offsets) for $\TTT$; the handover
        point moves to $x = Dr/(1+r)$, $r = 10^{H/(10n)}$.
  \item With 8\,dB shadowing a 3\,dB margin is crossed by chance 40\% of the time at the border:
        TTT and filtering are essential.
  \item X2 handover: report $\to$ request/ack $\to$ command $\to$ forwarding $\to$ RACH $\to$ complete
        $\to$ path switch $\to$ release.
  \item Too-late and ping-pong pull parameters in opposite directions; MRO, CHO and DAPS address
        the dilemma.
\end{itemize}
\end{takeawaybox}

% ---------------------------------------------------------------------
\section{Exercises}

\begin{exercisebox}{1}{Who decides?}\label{exr:who}
For each of the following, state whether the UE or the network takes the decision and on which
information: (a) LTE cell reselection; (b) LTE handover; (c) WCDMA active-set update; (d) 5G
conditional handover execution.
\end{exercisebox}

\begin{exercisebox}{1}{Two update procedures}\label{exr:lara}
Explain why a GPRS-capable GSM phone could perform both a location area update and a routing area
update at the same border, and how LTE avoided this duplication.
\end{exercisebox}

\begin{exercisebox}{2}{Offsets}\label{exr:offsets}
In \cref{ex:hopoint} the operator sets a cell-individual offset $O_\mathrm{cn} = +2$\,dB for the
neighbour (to off-load the serving cell). Where is the A3 condition now met? Which failure type
becomes more likely?
\end{exercisebox}

\begin{exercisebox}{2}{Filtering}\label{exr:filter}
Measurements are taken every 40\,ms with independent errors of standard deviation 4\,dB. Using
\cref{eq:l3filter}, compute the standard deviation after filtering for $K=4$ and $K=8$, and the
filter's time constant (number of samples for the impulse response to fall to $1/\e$).
\end{exercisebox}

\begin{exercisebox}{3}{Optimal area with two classes}\label{exr:twoclass}
A tracking area serves 70\% pedestrians ($v=1.5$\,m/s, 1 call/hour) and 30\% vehicles
($v=15$\,m/s, 3 calls/hour). Using \cref{eq:cost} with the other parameters of \cref{ex:la},
find the single $N$ that minimises the average cost, and the optimal $N$ if each class had its
own TA list. How much does per-class tailoring save?
\end{exercisebox}

\begin{exercisebox}{3}{Too late}\label{exr:late}
Assume RLF occurs if the serving RSRP falls 8\,dB below the neighbour before the handover command
is received, and that preparation plus command delivery take 80\,ms after the report. For the
geometry of \cref{ex:hopoint} without shadowing, find the maximum speed for which no too-late
failure occurs with $H = 3$\,dB and $\TTT = 320$\,ms. Discuss which parameter you would change for
a high-speed train line.
\end{exercisebox}

\subsection*{Solutions}
\addcontentsline{toc}{subsection}{Solutions}

\begin{solutionbox}{\ref{exr:who}}
(a) UE, using broadcast parameters ($Q_{\mathrm{rxlevmin}}$, $Q_{\mathrm{hyst}}$, priorities) and its own
measurements. (b) Source eNB, using UE measurement reports and its own load/admission knowledge.
(c) RNC, using UE event reports (1A/1B/1C). (d) The UE executes, but only towards targets prepared by
the network and when the network-configured condition holds -- a shared decision.
\end{solutionbox}

\begin{solutionbox}{\ref{exr:lara}}
LA and RA belonged to two separate core domains (MSC/VLR and SGSN) with separate procedures. Unless
the optional Gs interface allowed the SGSN to update the MSC on the phone's behalf (combined
updates), the phone signalled twice. LTE has a single packet core with one mobility entity (MME)
and one area concept (TA); CS-related registration, where needed for CS fallback, is performed by
the MME over SGs.
\end{solutionbox}

\begin{solutionbox}{\ref{exr:offsets}}
The condition becomes $M_\mathrm{n} + 2 - M_\mathrm{p} > 3$, i.e.\ $H_{\mathrm{eff}} = 1$\,dB:
$r = 10^{1/35} = 1.068$, $x_{\mathrm{HO}} = 1000\times1.068/2.068 = 516$\,m. Handovers happen earlier;
ping-pong and too-early failures become more likely, too-late less likely.
\end{solutionbox}

\begin{solutionbox}{\ref{exr:filter}}
Variance factor $a/(2-a)$. $K=4$: $a=0.5$, factor $1/3$, $\sigma_F = 4/\sqrt3 = 2.31$\,dB.
$K=8$: $a = 0.25$, factor $1/7$, $\sigma_F = 1.51$\,dB. Impulse response $(1-a)^k$ falls to
$1/\e$ after $k = -1/\ln(1-a)$: 1.44 samples (58\,ms) for $K=4$, 3.48 samples (139\,ms) for $K=8$.
Better smoothing costs reaction time.
\end{solutionbox}

\begin{solutionbox}{\ref{exr:twoclass}}
Per-user cost $C_i(N) = 4c_\mathrm{u}v_i/(\pi a\sqrt N) + c_\mathrm{p}\lambda_iN$. The average has the same
form with $\bar v = 0.7\times1.5+0.3\times15 = 5.55$\,m/s and $\bar\lambda = 0.7\times1+0.3\times3 =
1.6$ calls/h, so $N^\ast = (2\times5.55/(\pi\times500\times0.1\times1.6/3600))^{2/3} = 29.3$, cost
$\approx 14.1$ units per user-hour. Separately: pedestrians $N^\ast = 16.8$ (cost 5.0), vehicles
$N^\ast = 37.5$ (cost 33.7), average $0.7\times5.0+0.3\times33.7 = 13.6$. Tailoring saves only about
3\% here, because the cost curve is flat near its optimum; it pays off more when the classes differ
much more (e.g.\ stationary IoT devices versus trains).
(The script \texttt{scripts/validate\_examples.py} reproduces these values.)
\end{solutionbox}

\begin{solutionbox}{\ref{exr:late}}
The RLF point is where $M_\mathrm{n}-M_\mathrm{p} = 8$\,dB: $r = 10^{8/35} = 1.693$,
$x_{\mathrm{RLF}} = 1000\times1.693/2.693 = 629$\,m. The report is triggered at 549\,m and the command
arrives after $\TTT + 80$\,ms $= 0.40$\,s. The UE must not travel more than $629-549 = 80$\,m in that
time: $v_{\max} = 80/0.40 = 200$\,m/s. Without shadowing there is ample margin; with shadowing,
fading of the serving cell can cut this distance drastically. For a train line, one would reduce
TTT (speed-dependent scaling), lower the offset, deploy cells along the track with larger overlap,
or use conditional handover.
\end{solutionbox}

% ---------------------------------------------------------------------
\section{Further reading}
LTE mobility is described in \cite{ts36300}, idle-mode procedures in \cite{ts36304} and the
measurement events in \cite{ts36331}; GSM/UMTS handover in \cite{ts23009}; NR mobility, CHO and DAPS
in \cite{ts38300}. Akyildiz et al.\ \cite{akyildiz1999mobility} survey location management and the
fluid-flow model. Sesia et al.\ \cite{sesia2011lte} and Holma and Toskala \cite{holma2011lte}
discuss handover parameter tuning in LTE.

\printbibliography[heading=subbibliography,title={\ACNstrRefs}]
\end{document}
